% Part: incompleteness
% Chapter: representability-in-q
% Section: basic-representable

\documentclass[../../../include/open-logic-section]{subfiles}

\begin{document}

\olfileid{inc}{req}{bre}
\olsection{Fungsi Dhasar Representabel ing~$\Th{Q}$}

Luwih dhisik, kita kudu nuduhake yen kabeh fungsi dhasar representabel
ing~$\Th{Q}$. Pungkasane, kita kudu nuduhake carane masangake saben
fungsi dhasar $k$-argumen $f(x_0,\dots,x_{k-1})$ karo !!a{formula}
$!A_f(x_0,\dots,x_{k-1},y)$ kang makili fungsi iku.

Kita bakal bisa makili nol, penerus, panjumlahan,
perkalian, fungsi karakteristik kanggo kesamaan, lan proyeksi. Ing
saben kasus, formula kang makili fungsi kasebut cukup
langsung; umpamane, nol diwakili dening formula $y = \Obj 0$,
penerus dening !!{formula} $x_0' = y$, lan panjumlahan dening
!!{formula} $(x_0 + x_1) = y$. Gaweyan kang perlu yaiku nuduhake yen
$\Th{Q}$ bisa mbuktekake !!{sentence}s kang cocog; umpamane, nyatakake
yen panjumlahan diwakili dening !!{formula} ing ndhuwur
mbutuhake bukti yen kanggo saben pasangan wilangan natural $m$ lan $n$,
$\Th{Q}$ mbuktekake
\begin{align*}
& \eq[\num n + \num m][\num {n+m}] \text{ lan}\\
& \lforall[y][(\eq[(\num n + \num m)][y] \lif \eq[y][\num{n+m}])].
\end{align*}

\begin{prop}
\ollabel{prop:rep-zero}
Fungsi nol $\Zero(x) = 0$ diwakili ing~$\Th{Q}$ dening
$!A_{\Zero}(x,y) \ident \eq[y][\Obj 0]$.
\end{prop}

\begin{prop}
\ollabel{prop:rep-succ}
Fungsi penerus $\Succ(x) = x+1$ diwakili ing~$\Th{Q}$ dening
$!A_{\Succ}(x,y) \ident \eq[y][x']$.  
\end{prop}

\begin{prop}
\ollabel{prop:rep-proj}
Fungsi proyeksi $\Proj{n}{i}(x_0, \dots, x_{n-1}) = x_i$
diwakili ing~$\Th{Q}$ dening \[!A_{\Proj{n}{i}}(x_0, \dots, x_{n-1}, y) \ident \eq[y][x_i].\]
\end{prop}

\begin{prob}
Buktekna yen $\eq[y][\Obj 0]$, $\eq[y][x']$, lan $\eq[y][x_i]$
makili $\Zero$, $\Succ$, lan $\Proj{n}{i}$ miturut urutane.
\end{prob}

\begin{prop}
\ollabel{prop:rep-id}
Fungsi karakteristik saka~$=$,
\[
\Char{=}(x_0, x_1) =
\begin{cases}
  1 & \text{yen } x_0 =x_1\\
  0 & \text{yen ora}
\end{cases}
\]
diwakili ing~$\Th{Q}$ dening
\[
  !A_{\Char{=}}(x_0, x_1, y) \ident (\eq[x_0][x_1] \land \eq[y][\num{1}]) \lor (\eq/[x_0][x_1] \land
\eq[y][\num{0}]).
\]
\end{prop}

Buktine mbutuhake lema iki.

\begin{lem}
\ollabel{lem:q-proves-neq} Kanggo wilangan natural $n$ lan $m$, yen $n
\neq m$, mula $\Th{Q} \Proves \eq/[\num n][\num m]$.
\end{lem}

\begin{proof}
Gunakna induksi marang $n$ kanggo nuduhake yen kanggo saben $m$, yen
$n \neq m$, mula $Q \Proves \eq/[\num n][\num m]$.

Ing kasus dhasar, $n = 0$. Yen $m$ ora padha karo $0$, mula $m = k +
1$ kanggo sawijining wilangan natural $k$. Kita duwe aksioma kang nyatakake
$\lforall[x][\eq/[0][x']]$. Kanthi aksioma kuantor, yen $x$ diganti
  $\num k$, kita bisa nyimpulake $\eq/[0][\num k']$. Nanging $\num k'$
  iku mung $\num m$.

Ing langkah induksi, kita bisa nganggep pratelan iki bener kanggo $n$,
banjur nimbang $n+1$. Tetepna $m$ minangka wilangan natural sembarang. Ana rong
prakara kang bisa kelakon: $m = 0$ utawa ana $k$ saengga $m = k+1$. Kasus kapisan
ditangani kaya ing ndhuwur. Ing kasus kapindho, upamana $n+1 \neq
k+1$. Mula $n \neq k$. Miturut hipotesis induksi kanggo $n$, kita nduweni
$\Th{Q} \Proves \eq/[\num n][\num k]$. Kita duwe aksioma kang nyatakake
$\lforall[x][\lforall[y][\eq[x'][y'] \lif \eq[x][y]]]$. Kanthi
aksioma kuantor, kita entuk $\eq[\num n'][\num k'] \lif \eq[\num n][\num
  k]$. Kanthi logika proposisional, kita bisa nyimpulake, ing $\Th{Q}$,
$\eq/[\num n][\num k] \lif \eq/[\num n'][\num k']$. Kanthi modus
ponens, kita nyimpulake $\eq/[\num n'][\num k']$, kaya kang dikarepake,
amarga $\num k'$ iku $\num m$.
\end{proof}

\begin{explain}
Gatekna yen lema iki ora nyatakake akeh: pokoke, $\Th{Q}$ bisa
mbuktekake yen numeral kang beda nuduhake objek kang beda. Umpamane,
$\Th{Q}$ mbuktekake $0'' \neq 0'''$. Nanging kanggo nuduhake yen iki
bener kanthi umum perlu ngati-ati. Gatekna uga yen sanajan kita
nggunakake induksi, induksi iki ana \emph{ing njaba} $\Th{Q}$.
\end{explain}

\begin{proof}[Bukti \olref{prop:rep-id}]
Yen $n = m$, mula $\num{n}$ lan $\num{m}$ iku suku kang padha, lan
$\Char{=}(n, m) = 1$. Nanging $\Th{Q} \Proves (\eq[\num{n}][\num{m}] \land
\eq[\num{1}][\num{1}])$, mula teori iku mbuktekake $!A_=(\num{n}, \num{m},
\num{1})$. Yen $n \neq m$, mula $\Char=(n, m) = 0$. Miturut
\olref{lem:q-proves-neq}, $\Th{Q} \Proves \eq/[\num{n}][\num{m}]$ lan
uga $(\eq/[\num{n}][\num{m}] \land \Obj 0 = \Obj 0)$. Dadi $\Th{Q}
\Proves !A_=(\num{n}, \num{m}, \num{0})$.

Kanggo bagean kapindho, ana rong kasus maneh. Yen $n = m$, kita kudu
nuduhake yen $\Th{Q} \Proves \lforall[y][(!A_=(\num{n}, \num{m}, y)
  \lif \eq[y][\num{1}])]$. Kanthi argumentasi informal, upamana $!A_=(\num{n},
\num{m}, y)$, yaiku,
\[
(\eq[\num{n}][\num{n}] \land \eq[y][\num{1}]) \lor
(\eq/[\num{n}][\num{n}] \land \eq[y][\num{0}])
\]
Disjungsi sisih kiwa njalari $\eq[y][\num{1}]$ kanthi logika; sisih
tengen mbantah $\eq[\num{n}][\num{n}]$ kang bisa dibuktekake kanthi
logika.

Kosok baline, upamana $n \neq m$. Mula $!A_=(\num{n},
\num{m}, y)$ iku
\[
(\eq[\num{n}][\num{m}] \land \eq[y][\num{1}]) \lor
(\eq/[\num{n}][\num{m}] \land \eq[y][\num{0}])
\]
Ing kene, disjungsi sisih kiwa mbantah $\eq/[\num{n}][\num{m}]$, kang
bisa dibuktekake ing $\Th{Q}$ miturut \olref{lem:q-proves-neq};
disjungsi sisih tengen njalari $\eq[y][\num{0}]$.
\end{proof}

\begin{prop}
\ollabel{prop:rep-add}
Fungsi panjumlahan $\Add(x_0, x_1) = x_0+x_1$ diwakili
ing~$\Th{Q}$ dening
\[
  !A_{\Add}(x_0, x_1, y) \ident \eq[y][(x_0 + x_1)].
\]
\end{prop}

\begin{lem}
\ollabel{lem:q-proves-add}
$\Th{Q} \Proves \eq[(\num{n} + \num{m})][\num{n+m}]$
\end{lem}

\begin{proof}
Kita mbuktekake iki kanthi induksi marang~$m$. Yen $m = 0$, pratelane
yaiku $\Th{Q} \Proves \eq[(\num{n} + \Obj 0)][\num{n}]$. Iki lumaku
miturut aksioma~$!Q_4$. Saiki upamana pratelan iki bener kanggo $m$;
ayo mbuktekake kanggo $m+1$, yaiku $\Th{Q} \Proves \eq[(\num{n} +
  \num{m+1})][\num{n+m+1}]$. Gatekna yen $\num{m+1}$ iku mung $\num{m}'$,
lan $\num{n+m+1}$ iku mung $\num{n+m}'$. Miturut aksioma $!Q_5$, $\Th{Q}
\Proves \eq[(\num{n} + \num{m}')][(\num{n}+\num{m})']$. Miturut
hipotesis induksi, $\Th{Q} \Proves \eq[(\num{n} + \num{m})][\num{n+m}]$.
Mula $\Th{Q} \Proves \eq[(\num{n} + \num{m}')][\num{n+m}']$.
\end{proof}

\begin{proof}[Bukti \olref{prop:rep-add}]
!!{formula}~$!A_\Add(x_0, x_1, y)$ kang makili $\Add$ yaiku
$\eq[y][(x_0 + x_1)]$. Luwih dhisik, kita nuduhake yen $\Add(n, m) = k$, mula
$\Th{Q} \Proves !A_\Add(\num{n}, \num{m}, \num{k})$, yaiku $\Th{Q}
\Proves \eq[\num{k}][(\num{n} + \num{m})]$. Nanging amarga $k = n + m$,
$\num{k}$ iku mung $\num{n+m}$, lan kita wis nuduhake ing
\olref{lem:q-proves-add} yen $\Th{Q} \Proves \eq[(\num{n} +
  \num{m})][\num{n+m}]$.

Kita uga kudu nuduhake yen $\Add(n, m) = k$, mula
\[
\Th{Q} \Proves \lforall[y][(!A_\Add(\num{n}, \num{m}, y) \lif
  \eq[y][\num{k}])].
\]
Upamana kita nduweni $\eq[(\num{n} + \num{m})][y]$. Amarga
\[
\Th{Q} \Proves \eq[(\num{n}+\num{m})][\num{n+m}],
\]
kita bisa ngganti sisih kiwa karo $\num{n+m}$ lan entuk
$\eq[\num{n+m}][y]$, kanggo~$y$ sembarang.
\end{proof}

\begin{prop}
\ollabel{prop:rep-mult}
Fungsi perkalian $\Mult(x_0, x_1) = x_0 \cdot x_1$ diwakili
ing~$\Th{Q}$ dening
\[
  !A_{\Mult}(x_0, x_1, y) \ident y = (x_0 \times x_1).
\]
\end{prop}

\begin{proof} Gladhen. \end{proof}

\begin{lem}
\ollabel{lem:q-proves-mult}
$\Th{Q} \Proves \eq[(\num{n} \times \num{m})][\num{n \cdot m}]$
\end{lem}

\begin{proof} Gladhen. \end{proof}

\begin{prob}
Buktekna \olref[inc][req][bre]{lem:q-proves-mult}.
\end{prob}

\begin{prob}
Gunakna \olref[inc][req][bre]{lem:q-proves-mult} kanggo mbuktekake
\olref[inc][req][bre]{prop:rep-mult}.
\end{prob}

\begin{explain}
  Elinga yen kita nggunakake $\times$ minangka simbol fungsi ing basa
  aritmetika, lan $\cdot$ kanggo operasi perkalian biasa marang
  wilangan. Dadi $\cdot$ bisa katon ing antarane ekspresi wilangan
  (kayata ing $m \cdot n$), dene $\times$ mung katon ing antarane suku
  basa aritmetika (kayata ing $(\num{m} \times
  \num{n})$). Sing luwih mbingungake, $+$ digunakake kanggo
  !!{function} lan uga operasi panjumlahan. Yen katon ing antarane
  suku---umpamane, ing $(\num{n} + \num{m})$---iku !!{function} kanthi
  $2$ argumen ing basa aritmetika, lan yen katon ing antarane
  wilangan---umpamane, ing $n+m$---iku operasi panjumlahan.
  Iki kalebu kasus $\num{n+m}$: iki numeral baku
  kang cocog karo wilangan~$n+m$.
\end{explain}

\end{document}
